\documentclass{article}
\usepackage[T1]{fontenc}
\usepackage{lmodern}
\usepackage{graphicx}
\usepackage{amsmath,amssymb}
\usepackage[a4paper,margin=2cm]{geometry}
\usepackage[absolute,overlay]{textpos}
\setlength{\TPHorizModule}{1mm}
\setlength{\TPVertModule}{1mm}
\usepackage[indonesian]{babel}
\usepackage{hyperref}
\setlength{\emergencystretch}{2em}
% unit: Halaman 9

\begin{document}

% === Halaman 9 (llm) ===

\begin{itemize}
    \item Dengan panjang kunci 128-bit, maka terdapat sebanyak
    \[ 2^{128} = 3,4 \times 10^{38} \text{ kemungkinan kunci.} \]
    \item Jika komputer tercepat dapat mencoba 1 juta kunci setiap detik, maka akan dibutuhkan waktu $5,4 \times 10^{24}$ tahun untuk mencoba seluruh kunci.
    \item Jika komputer tercepat yang dapat mencoba 1 juta kunci setiap milidetik, maka dibutuhkan waktu $5,4 \times 10^{18}$ tahun untuk mencoba seluruh kunci.
    \item Artinya, AES diharapkan tahan terhadap serangan \textit{exhaustive key search attack (brute force attack)}
\end{itemize}

\end{document}
